\documentclass[11pt]{article}
\usepackage[margin=1in]{geometry}
\usepackage{amsmath,amssymb}
\title{A Literal-Set Contradiction in Conjecture 00000007838}
\author{earthking11}
\date{October 7, 2026}
\begin{document}
\maketitle

\section*{The stated object}
The conjecture calls the displayed set
\[
  H_{\mathrm{star}}=\{0,2,4,\ldots\}
\]
a hyperelliptic Weierstrass semigroup and asserts that it is realized at every point. With the ordinary meaning of the ellipsis, this is the set of all nonnegative even integers. The English and Chinese statements use the same displayed set.

\section*{Why it cannot be a Weierstrass semigroup}
A numerical semigroup is an additive submonoid of $\mathbb N_0$ with finite complement. The Weierstrass semigroup at a point of a smooth curve is a numerical semigroup; for a genus-$g$ curve its complement has exactly $g$ elements, by the Weierstrass gap theorem. Thus finite complement is a necessary condition, independent of the formula conjectured for $R(H)$.

For every $k\in\mathbb N_0$, the number $2k+1$ is odd and hence is not in $H_{\mathrm{star}}$. The map $k\mapsto 2k+1$ is injective, so there are infinitely many distinct missing integers. Equivalently,
\[
  \mathbb N_0\setminus H_{\mathrm{star}}
  \supseteq \{2k+1:k\in\mathbb N_0\},
\]
and the right-hand set is infinite. Therefore $H_{\mathrm{star}}$ is not cofinite and is not a numerical semigroup, hence cannot be a Weierstrass semigroup of any finite genus.

The conjecture's stated hyperelliptic example therefore cannot exist with the displayed set. Since this is one of the conjecture's asserted clauses, its full conjunction is false as written; no claim about the other clauses is needed.

\section*{Scope of the formalization}
The Lean project defines the displayed set as all even natural numbers, defines the necessary finite-gap condition, and proves that no Weierstrass semigroup can equal that set. This is the semantic bridge used for the disproof. If the intended object was instead a finite-genus hyperelliptic Weierstrass semigroup with a finite list of odd gaps and an eventual even tail, that is a different set from $\{0,2,4,\ldots\}$ and would require a corrected conjecture. This submission does not silently substitute that corrected object.

\section*{Reference for the standard definition}
E. Mart\'inez-Moro, \emph{Numerical Semigroups and Codes}, Chapter 5 of \emph{Algebraic Geometry Modeling in Information Theory}, World Scientific (2013), defines a numerical semigroup as a subset of $\mathbb N_0$ containing $0$, closed under addition, with finite complement, and identifies Weierstrass semigroups as examples. The Weierstrass gap theorem gives exactly $g$ gaps for a point on a genus-$g$ curve. A second account is N. Pflueger, \emph{Weierstrass Semigroups on Castelnuovo Curves}, Introduction, which states that the complement of a Weierstrass semigroup has size $g$.

\end{document}
